\section{Higher-dimensional extensions: available energy and momentum transfer}
\label{sec:3d}

Unbinned OmniFold adds an observable by adding a column to the step-1/step-2
classifiers, without an explicit migration matrix or per-bin response tensor.
It avoids the combinatorial binning penalty of matrix unfolding but is not free
of a dimensional cost: classifier training, acceptance/support coverage, sample
sparsity and uncertainty quantification become more demanding as observables
are added. The absence of a practical D'Agostini
IBU analog is the primary methodological motivation.
We extend the measurement to
$d^{3}\sigma/(d\pT\,d\pPar\,d\Eavail)$ with available energy as a third axis.

MINERvA has published binned triple-differential results
(Fig.~\ref{fig:landscape}, Table~\ref{tab:landscape}), including recent comparisons
between low- and medium-energy beams~\cite{MINERvA:2026qeleme}. The latter measure
$\pT$, $\pPar$, and available recoil energy for a quasielastic-like signal, for
which the available energy is the sum of proton kinetic energies. This differs
from the inclusive signal and species-dependent energy definition used here.
Our distinction is the unbinned, simultaneous CC-inclusive scalar formulation
extended to three, four and five recorded observables, not multidimensionality
by itself.

\begin{figure}[H]
  \centering
  \includegraphics[width=0.97\textwidth]{minerva_unfolding_landscape}
  \caption{Number of kinematic variables unfolded \emph{simultaneously} in
  surveyed MINERvA differential cross-section measurements, by year, colored by channel
  (a double-differential $d^{2}\sigma/dx\,dy$ counts as two).
  Prior results use explicit binned unfolding, including a triple-differential
  quasielastic-like measurement; this work extends an unbinned CC-inclusive
  scalar unfold to five recorded observables.
  Source: the MINERvA data-release page and the cited publications. The plotted
  inventory runs through 2025; newer measurements are listed in
  App.~\ref{app:landscape}. Classification uses each analysis's headline
  differential result; flux, reconstruction, generator-tuning and
  total-only entries are excluded.}
  \label{fig:landscape}
\end{figure}

\subsection{Operational definition and provenance}
\label{sec:3d-method}
\label{sec:eavail-definition}
Truth available energy follows the code; ``low-recoil convention'' is
shorthand. In MeV, \texttt{CVUniverse::GetEAvailableTrue()} implements
\begin{equation}
 E_{\rm avail}^{\rm true}=
 \sum_{p}(E-938.27)
 +\sum_{\pi^\pm}(E-135)
 +\sum_{\pi^0}E
 +\sum_{\gamma}E .
 \label{eq:eavail-implemented}
\end{equation}
Protons and charged pions contribute the indicated kinetic-energy-like terms;
neutral pions and photons contribute total energy. All other final-state species
contribute zero, including neutrons, the primary muon, secondary $e^\pm$,
neutrinos, kaons, antiprotons, strange baryons, other mesons and nuclear-remnant
records. Reconstructed available energy
is the MAT calorimetric quantity
\begin{equation}
 E_{\rm avail}^{\rm reco}=1.17\left[(E_{\rm blob}^{\rm tracker}
 -E_{\mu\text{-fuzz}}^{\rm tracker})+(E_{\rm blob}^{\rm ECAL}
 -E_{\mu\text{-fuzz}}^{\rm ECAL})\right],
 \label{eq:eavail-reco-implemented}
\end{equation}
implemented by \texttt{CVUniverse::NewEavail()}; the code applies no final
non-negative clamp.

Both functions were inherited from MAT-MINERvA rather than selected in this
analysis.  Agreement with a published convention is therefore a defence of the
inherited position, not its provenance, and there is no unique published
$E_{\rm avail}$ definition.  Rodrigues \emph{et al.}~\cite{MINERvA:2016lowrecoil}
give a closed five-species list of proton and charged-pion kinetic energy plus
$\pi^0$, electron and photon total energy.  Equation~\eqref{eq:eavail-implemented}
matches four of those five species and excludes secondary $e^\pm$.
The later Ascencio \emph{et al.}~\cite{MINERvA:2022incl} definition is open and
includes other final-state particles except neutrons, subtracting a nucleon
mass for strange baryons. It therefore includes species our closed list omits.
The charged-pion term in Rodrigues and in Henry/Su \emph{et al.}
Eq.~4~\cite{MINERvA:2023avz} is kinetic energy; the
\texttt{minerva-ml} implementation discussed in \S\ref{sec:pet-gregor} instead
uses charged-pion total energy.  We therefore do not claim generically to
``match the published definition.''

Production subtracts the inherited constant
\SI{135}{MeV} per charged pion rather than the physical charged-pion mass of
approximately \SI{139.57}{MeV}; this was not corrected. The effect is measured
but has not been assessed against the adopted scalar-5D projection. Its
materiality is not established, and this note does not call it immaterial.
% This note neither applies the six-site correction nor triggers a rerun.

\paragraph{Pipeline and binning.}
The event loop evaluates both accessors over all twelve ME~FHC playlists,
carrying $\Eavail$ alongside $\pT$/$\pPar$. The 2D driver uses a three-column
feature stack $(\pT,\pPar,\Eavail)$; extraction uses the 3D Jacobian
$\Delta\pT\,\Delta\pPar\,\Delta\Eavail$. Flux $\Phi$ remains applied per
$\pT$ (broadcast over $\pPar$ and $\Eavail$). The $\Eavail$ binning
$[0,0.1,0.2,0.4,0.8,1.5,3.0,100]\,\si{GeV}$ ends in a wide catch bin to
capture the full CC-inclusive recoil tail in the $\Eavail$-marginal. Global
OmniFold completeness is $c=1.0000$, as in 2D.

\subsection{The available-energy spectrum}
\label{sec:3d-spectrum}
The unfolded $d\sigma/d\Eavail$, the 1D projection
of the 3D result
(Fig.~\ref{fig:eavail}), falls monotonically from
\SI{2.19e-38}{cm^2/GeV/nucleon} at the low-recoil peak to
\SI{6.8e-41}{} in the catch bin, the expected CC-inclusive shape with most
of the rate at $\Eavail\lesssim\SI{0.8}{GeV}$.

\begin{figure}[H]
  \centering
  \includegraphics[width=0.66\textwidth]{eavail_spectrum}
  \caption{Unfolded $d\sigma/d\Eavail$, the 1D projection of the 3D measurement.
  The wide $[3,100]\,\si{GeV}$ catch bin is excluded from the axis; its value is annotated.}
  \label{fig:eavail}
\end{figure}

\subsection{Validation against the 2D measurement}
\label{sec:3d-anchor}
Marginalizing over $\Eavail$ must reproduce the 2D result of
\S\ref{sec:results}, our validation anchor in the absence of a directly matched
published CC-inclusive 3D reference for this signal and phase space.
The \emph{normalization} agrees closely, with total $\sigma$ at
$+\SI{0.95}{\percent}$ and median per-bin marginal/2D ratio $1.0016$.
The \emph{shape} agrees less closely. Using published covariance gives an indicative distance
$\chi^{2}/\mathrm{ndf}=4.98$,
above the frozen 2D values ($3.66$ default, $2.65$ LightGBM-CV), driven by
$\sim\SI{4.4}{\percent}$ per-bin scatter introduced by the third axis in the
$(\pT,\pPar)$ marginal (Fig.~\ref{fig:eavail-marg}).
Shared data and unavailable cross-covariance prevent interpreting this number
as a calibrated goodness-of-fit statistic.

\begin{figure}[H]
  \centering
  \includegraphics[width=0.7\textwidth]{eavail_marginal_vs_paper_pull_full}
  \caption{Standardized per-bin differences of the $\Eavail$-marginal, scaled by
  published uncertainties, on the $14\times16$ grid and in distribution.
  Shared data and unavailable cross-covariance make these descriptive;
  they are not formal independent-result pulls.}
  \label{fig:eavail-marg}
\end{figure}

\subsection{Closure: the one injected $\Eavail$ deformation is recovered}
\label{sec:3d-closure}
A reweight closure isolates whether the $\SI{4.4}{\percent}$ scatter reflects
method bias or real information. The 3D unfold recovers a known
$+\SI{30}{\percent}$ Gaussian bump injected into pseudo-data in \emph{truth}
$\Eavail$ (center \SI{0.3}{GeV}, width \SI{0.15}{GeV}).
The recovered step-2 mean weight is $1.048$ against the injected $1.0485$.
Residuals against reweighted truth on the $\Eavail$-marginal have median
$0.9999$ and standard deviation \textbf{\num{0.0006}}
(\SI{0.06}{\percent}); the per-3D-bin residual standard deviation is
\SI{3.2}{\percent} from MC statistics in sparse high-recoil corners.
Marginal scatter is $\sim70\times$ tighter than the
$\SI{4.4}{\percent}$ data-versus-2D scatter; total normalization matches
to better than \SI{1}{\percent}. In the metric of \S\ref{sec:validation}
($\Eavail$ projection) the residual is \SI{1.84}{\percent} vs \SI{18.1}{\percent}
injected (ratio $0.102$); the 3D test passes.

\emph{What closes, stated exactly.} The deformation recovered is one shape: a
$+\SI{30}{\percent}$ Gaussian in truth $\Eavail$ at
center~\SI{0.3}{GeV}, width~\SI{0.15}{GeV}, injected into the pseudo-data and
unfolded \emph{with the same detector response the unfold was given}. On that
test the method does not manufacture or destroy the feature, which is what
makes the data-versus-2D scatter readable as
data\,$\leftrightarrow$\,MC structure resolved by the available-energy
dimension rather than as an artifact of this particular injection. It does
\emph{not} exclude estimator bias for other deformations or test response
mismatch, and the hidden-variable closure exceeds its threshold at the largest
amplitude (\S\ref{sec:validation}, App.~\ref{app:response-mismatch}).

\subsection{Classifier-family cross-check (GBDT vs.\ NN)}
\label{sec:3d-classifier}
Our gradient-boosted decision trees differ from the neural-network classifiers
used in published OmniFold analyses~\cite{Andreassen:2019cjw,unbinnedguide}.
GBDT output must be a calibrated posterior for $w=p/(1-p)$ to hold.
The step-1 reco classifier (data vs.\ MC reco) has a near-diagonal reliability
curve, Brier \(\approx0.25\) at AUC \(\approx0.537\).
Data and reco MC are nearly indistinguishable, so the reweight is a modest correction.
A full scalar Keras-MLP cross-check with the same weighted 3D inputs and
five-iteration OmniFold loop gives total ratio $1.0078$ versus GBDT and median projection
differences of \SIrange{0.7}{1.4}{\percent} (Fig.~\ref{fig:calibration}).
Agreement in the populated bulk establishes that the scalar result is not
specific to the production classifier family. Section~\ref{sec:pet-shape}
reports the diagnostic cross-check using low-level recoil point clouds in
place of engineered scalars.

\begin{figure}[H]
  \centering
  \includegraphics[width=0.95\textwidth]{nn_vs_gbdt_full}
  \caption{Five-iteration scalar OmniFold cross-check with identical weighted
  3D inputs and extraction, changing only the classifier family. Top:
  absolutely normalized GBDT and Keras-MLP projections; bottom: binwise ratio.
  The MLP/GBDT total-cross-section ratio is $1.0078$; median
  absolute projection differences are \SI{1.20}{\percent} in $\pT$,
  \SI{1.36}{\percent} in $\pPar$, and \SI{0.66}{\percent} in $\Eavail$.
  Larger excursions occur only in sparsely populated outermost
  $\pT$ and $\pPar$ bins.}
  \label{fig:calibration}
\end{figure}

\subsection{Generator comparison in 3D}
\label{sec:3d-models}
The ancillary provides only binned 2D model histograms. Quantitative 3D
comparisons or arbitrary projections require generating truth-level
events, histogrammed as $d^{3}\sigma$ in the same fiducial phase space with the
same truth $\Eavail$ definition. The unfolded result is at truth level, so comparisons need no detector
simulation. We generated and processed GENIE, NuWro
and GiBUU samples, and extracted MINERvA Tune~v1 from analysis MC truth with
its full tune weight. No spectra were digitized or externally supplied.
\begin{itemize}
  \item \textbf{GENIE 2.12.10 central value}, generated with \texttt{gevgen}
        on a CH target with the MINERvA ME~FHC flux, using the
        DefaultPlusValenciaMEC cross-section splines~\cite{Andreopoulos:2009rq}
        but gevgen's default event-generator list, which excludes MEC, so
        the sample has no 2p2h events (\S\ref{sec:3d-2p2h}); it has no MINERvA reweights;
        events are reweighted to the analysis flux as described below.
  \item \textbf{MINERvA Tune v1}, the analysis tune (base GENIE $+$ RPA,
        low-recoil 2p2h, and non-resonant-pion reweights), extracted directly
        from event-loop MC truth, which
        % (the \texttt{w\_truth} weights)
        reproduces the shipped ancillary curve to \SI{0.01}{\percent} in
        normalization.
  \item \textbf{NuWro 21.09}~\cite{Golan:2012wx}, an independent generator
        on C with the same flux (after reweighting below) and phase space.
  \item \textbf{GiBUU 2019}~\cite{Buss:2011mx}, a transport-theory generator
        on C, built and run \emph{natively} on the compute system without a
        container, with the same flux and phase space.
\end{itemize}
NEUT~\cite{Hayato:2009zz} is not included.
Each prediction is normalized by its own flux-averaged total CC cross section
per nucleon.

\paragraph{Flux repair of the GENIE and NuWro predictions.}
% (2026-09-27)
GENIE (CV and $+$Valencia~2p2h) and NuWro were initially generated from a flux
without PPFX weights, capped at \SI{50}{GeV}. Variable-width bins stored flux
\emph{densities} but were sampled and normalized by \emph{content}.
Under-weighting every wide high-energy bin suppressed high-$\pPar$ predictions
and integrated rates (App.~\ref{app:provenance}, \texttt{flux-repair}).
% KNOWN_ISSUES.md row 83
GENIE and NuWro predictions throughout \S\S\ref{sec:3d}--\ref{sec:eavailw} are
repaired by exact per-event reweighting to the analysis PPFX central-value flux
on 0--\SI{100}{GeV}, with correctly sampled 50--\SI{100}{GeV} supplements.
The FSI-dial studies, departure ratio of \S\ref{sec:val-negweight} and
NuWro-shaped prior of \S\ref{sec:fps} retain unrepaired samples, as stated there.
GiBUU reads its own flux and receives an energy-dependent per-event reweight
to the same convention, changing its integrated cross section by a factor
$0.9966$. Its sample still stops at $E_\nu\approx\SI{20}{GeV}$, so it is empty
above $\pPar\approx\SI{20}{GeV/c}$.
% KNOWN_ISSUES.md row 82
Comparisons here and in \S\ref{sec:eavailw-band} were rebuilt with unchanged
figure producers on repaired inputs (App.~\ref{app:history}).
% rebuilt on 2026-09-27; ledger rows VL156--VL160 (App. E flux-repair), which supersede VL35--VL38 (App. H, pre-repair values)
The analysis simulation, MINERvA Tune~v1, is unaffected.

Figure~\ref{fig:3dmodels} compares central values only. Totals over the full
measurement phase space follow
GiBUU ($\SI{2.22e-38}{}$) $<$ NuWro ($\SI{2.66e-38}{}$) $<$ Tune v1
($\SI{2.71e-38}{}$) $<$ GENIE CV ($\SI{2.76e-38}{}$) $<$ data
($\SI{3.08e-38}{cm^2/nucleon}$). All four under-predict the rate by
\SIrange{10}{28}{\percent} of the data. Integrating the reported $\Eavail$ range
with the catch bin dropped gives the same order. All closely track the
$(\pT,\pPar)$ shape. The $\Eavail$ axis discriminates among them;
NuWro and Tune v1 \emph{suppress} the lowest-$\Eavail$ (quasielastic) bin
below bare GENIE through RPA and nuclear effects, while GiBUU does not.
Data lie \emph{above} all four across the quasielastic--$\Delta$ dip,
$0.1\le\Eavail<\SI{0.4}{GeV}$ (in the lowest bin GiBUU reaches the data).
The $\Eavail$ axis exposes this excess, invisible in 2D.

\begin{figure}[H]
  \centering
  \includegraphics[width=0.99\textwidth]{generators_vs_unfolded_band_central}
  \caption{1D projections of the unfolded 3D central cross section (black)
  versus truth-level predictions from GENIE 2.12 CV, MINERvA Tune v1, NuWro 21.09, and
  GiBUU 2019 (GENIE and NuWro: the flux-repaired predictions,
  \S\ref{sec:3d-models}). Models track the $\pT$ and $\pPar$ shapes but differ
  on $\Eavail$; all under-predict the quasielastic--$\Delta$
  dip ($0.1\le\Eavail<\SI{0.4}{GeV}$). The wide
  $\Eavail$ catch bin is omitted from the axis.
  Data have no error bars because the quotable
  3D covariance is not yet adopted (App.~\ref{app:history}).}
  \label{fig:3dmodels}
\end{figure}

For the integrated $\Eavail$ rate without the catch bin, GENIE CV, Tune~v1,
NuWro and GiBUU central values lie respectively
$\SI{7.2}{\percent}$, $\SI{9.5}{\percent}$, $\SI{12.8}{\percent}$ and
$\SI{22.2}{\percent}$ below the data.
% (flux-repaired predictions; ledger row \texttt{VL159})
No 3D covariance, generator $\chi^{2}$, $p$-value or significance is quoted:
the quotable 3D covariance is the exact projection of the adopted 5D covariance,
and that projection is constructed but not adopted (App.~\ref{app:history}).

\subsection{The low-available-energy excess recovers the known low-recoil 2p2h deficit}
\label{sec:3d-2p2h}
The low-$\Eavail$ excess is the most physically interesting feature the third
axis exposes, and the 3D framework lets us test candidate explanations by
generating the relevant physics as well as reweighting. MINERvA first isolated the rate enhancement between quasielastic
and $\Delta$ peaks, with multi-nucleon final states, in a
low-three-momentum-transfer subsample~\cite{MINERvA:2016lowrecoil}.
The empirical tune was built to absorb it; our muon-kinematics
$\times\,\Eavail$ measurement re-exposes that established low-recoil feature.

\paragraph{The tested pion-FSI variations do not account for it.}
A \texttt{grwght1p} reweight of 2M central-value GENIE events gives sub-percent
changes in $d\sigma/d\Eavail$ for both tested pion-FSI dials.
\texttt{FrInel\_pi} shifts the spectrum by
$\le\SI{0.74}{\percent}$ (peak in $0.10$--$\SI{0.20}{GeV}$) and \texttt{FrAbs\_pi}
by $\le\SI{0.82}{\percent}$ (peak in the lowest bin). Both use the pre-repair
GENIE sample (\S\ref{sec:3d-models}) as same-sample ratios, where the flux
defect cancels to first order; neither was recomputed. The integrated rate
% flux repair: KNOWN_ISSUES.md row 83
moves only $\pm\SI{0.02}{}$--$\SI{0.03}{\percent}$ at $\pm1\sigma$.
\footnote{\texttt{FrInel\_pi} is commented out of the standard MAT-MINERvA systematics
registry. In 2026 the collaboration confirmed this remains current practice
and that the reason is a degeneracy (overlapping pion-FSI dials require
dropping one), rather than a defect specific to that knob.}
A less-than-\SI{1}{\percent} knob cannot bridge an excess of
\SIrange{8}{36}{\percent} per bin over the GENIE central value below
$\Eavail=\SI{0.4}{GeV}$. The tested variations do not account for the excess,
a conclusion limited to these dials and their range. Other transport models
or proton-FSI effects are not excluded. Recent MINERvA quasielastic-like
measurements point toward overestimated proton and charged-pion FSI in their model
comparisons~\cite{MINERvA:2026qeleme}. Their different signal definition prevents
directly identifying that effect with our inclusive discrepancy.

\paragraph{The shape and size match 2p2h.} The base GENIE central value uses
the default event-generator list and contains \emph{no}
meson-exchange-current (2p2h/MEC) events
($\texttt{mec}=0$ for all $1.48\times10^{6}$ CC events of the generated
sample, whose unweighted modes, before the flux reweight of
\S\ref{sec:3d-models}, are QE
\SI{11}{\percent}, RES \SI{22}{\percent}, DIS \SI{66}{\percent}, COH
\SI{0.9}{\percent}; the 50--\SI{100}{GeV} supplement is MEC-free as well).
2p2h must therefore be \emph{added}; reweighting cannot supply it.
The mode decomposition of central-value $d\sigma/d\Eavail$ (per-bin mode
count-fractions $\times$ bin value), overlaid with data, locates the largest
positive gaps in the quasielastic--$\Delta$ dip,
$[0.10,0.20)$ and $[0.20,0.40)$, while central values are much closer in the
QE-dominated $[0,0.10)$ and $[0.40,1.50)$ regions.
At central value, \SI{70}{\percent} of the positive deficit lies at
$\Eavail\le\SI{0.4}{GeV}$, exactly where 2p2h lives. Closing the integrated deficit
requires 2p2h at $\approx\SI{45}{\percent}$ of the QE rate
(\SIrange{68}{71}{\percent} in the two dip bins), the standard
empirical/Valencia-MEC magnitude, compared with sub-percent FSI shifts
(Fig.~\ref{fig:modedecomp}).

\paragraph{Enabling Valencia 2p2h improves the description.} We regenerated the GENIE central
value with empirical 2p2h enabled (\texttt{--event-generator-list
Default+CCMEC}, 2M events). Because the Nieves--Simo--Vacas MEC channel, though
present in the splines, is absent from the \texttt{gspl2root} total-CC graph, we
normalize by the non-MEC CC count, anchoring QE$+$RES$+$DIS$+$COH to the known total and
adding MEC through $1/(1-f_{\mathrm{MEC}})$ with
$f_{\mathrm{MEC}}=\SI{2.87}{\percent}$.
2p2h falls in the dip and brings all three low-$\Eavail$ bins toward data
(Fig.~\ref{fig:mec}), filling
\SI{52}{\percent} of the data$-$CV gap at $\Eavail\le\SI{0.4}{GeV}$ and \SI{63}{\percent} of the
integrated deficit (CV $-\SI{7.2}{\percent}\to-\SI{2.6}{\percent}$;
App.~\ref{app:history}). The prediction
% ledger row \texttt{VL160}
still falls short in the dip, especially $[0.10,0.20)$, but its integral over
the reported $\Eavail$ range is closer to data than
Tune~v1 ($-\SI{2.6}{\percent}$ against $-\SI{9.5}{\percent}$). Tune v1
layers MINERvA's empirical low-recoil 2p2h \emph{enhancement} and RPA on top of
the stock Valencia MEC --- i.e.\ stock Valencia 2p2h is real but under-strength
for MINERvA low recoil, which is precisely the known MINERvA-tune story and is
consistent with NuWro and GiBUU (native, un-enhanced 2p2h) sitting low in the
dip as well.

2p2h leaves the high-$\Eavail$ $[1.50,3.00)$ bin unchanged. There, data
exceed the GENIE CV by about \SI{9}{\percent} (data/CV $1.09$;
App.~\ref{app:history}) and exceed Tune~v1 as well; the high-$\Eavail$
excess over Tune~v1 is localized on the $(\Eavail,W)$ plane in
\S\ref{sec:eavailw}.

The improvement supports the tested multinucleon contribution's relevance
without uniquely identifying a nuclear mechanism. For example, recent NEUT studies include
two-body currents contributing to one-particle--one-hole final states and examine their
interplay with axial form factors
\cite{FrancoMunoz:2026currents}. A two-body current is not synonymous with the
generator's 2p2h event category. Visible-energy comparisons can depend on the
primary interaction and subsequent hadronic transport.

Ascencio \emph{et al.}~\cite{MINERvA:2022incl} report the same deficit in
MINERvA's dedicated low-recoil inclusive $d^{2}\sigma/(dq_{3}\,d\Eavail)$
measurement for $q_{3}<\SI{1.2}{GeV}$ at the same
$\langle E_\nu\rangle\sim\SI{6}{GeV}$, finding GENIE and NuWro underpredicting
their low-$\Eavail$ data without enhanced 2p2h. Our muon-kinematics extension integrates all $q_3$
instead of fixing $q_3$, limiting comparison to $\Eavail$ \emph{shape} and
direction, without absolute normalization. The consistency confirms recovery
of the established MINERvA low-recoil/2p2h feature as a by-product of the
OmniFold 3D muon-kinematics unfolding. Bin-identical central values follow in
\S\ref{sec:4d}.

\begin{figure}[H]
  \centering
  \begin{minipage}{0.49\textwidth}
    \centering
    \includegraphics[width=\textwidth]{mode_decomp_eavail_central}
    \caption{Central-value $d\sigma/d\Eavail$ by interaction mode, overlaid
    with unfolded data. The excess lies in the quasielastic--$\Delta$ dip
    ($\Eavail\le\SI{0.4}{GeV}$), matching a missing 2p2h component's shape and
    size. Data carry no error bars because the quotable 3D covariance is not
    yet adopted (App.~\ref{app:history}).}
    \label{fig:modedecomp}
  \end{minipage}\hfill
  \begin{minipage}{0.49\textwidth}
    \centering
    \includegraphics[width=\textwidth]{compare_mec_eavail_central}
    \caption{GENIE regenerated with empirical Valencia 2p2h (red, stacked on
    no-2p2h CV; both flux-repaired) fills about half the low-$\Eavail$
    data$-$CV gap (\SI{52}{\percent}). The remainder occupies the region the
    MINERvA tune's empirical enhancement is built to absorb. Data carry no
    error bars because the quotable 3D covariance is not yet adopted
    (App.~\ref{app:history}).}
    \label{fig:mec}
  \end{minipage}
\end{figure}

\subsection{Fourth axis: three-momentum transfer $q_{3}$}
\label{sec:4d}
The fourth axis adds one feature (\S\ref{sec:3d}), three-momentum transfer
$q_{3}=|\vec{q}\,|$, chosen because its $(\Eavail,q_{3})$ plane separates
quasielastic, 2p2h and resonance processes and permits bin-compatible comparison to MINERvA's
dedicated low-recoil $d^{2}\sigma/dq_{3}\,d\Eavail$ measurement
(Ascencio \emph{et al.}~\cite{MINERvA:2022incl}). Truth $q_{3}$ comes from
the generator (\texttt{CVUniverse::Getq3True()}); reco $q_{3}$ uses muon
$Q^{2}$ and calorimetric recoil (\texttt{CVUniverse::RecoQ3()}). Both accessors
run over all twelve playlists, as for $\Eavail$. Lateral systematic shifts
change $q_{3}$ but leave $\Eavail$ invariant, requiring per-universe
shifted-$q_{3}$ dumps instead of central-value reuse.
The feature stack is $(\pT,\pPar,\Eavail,q_{3})$, with deliberately coarse
$q_{3}$ binning $[0,0.2,0.4,0.6,0.8,1.2,2.0,100]\,\si{GeV}$
(seven bins including a final catch bin), giving
$14\times16\times7\times7=10\,976$ four-dimensional bins (4830 populated).

\paragraph{Marginal anchor (4D$\to$3D).} The $q_{3}$ marginalization reproduces
the 3D $\Eavail$ spectrum and low-recoil excess (Fig.~\ref{fig:q3excess}, right).
The $\Eavail$-marginal data/Tune~v1 ratio is
$1.21$ in the lowest bin, dips to $1.01$ across the quasielastic--$\Delta$
region, and rises to $1.16$--$1.23$ in the high-$\Eavail$ DIS tail
(integrated data/Tune~v1 $=1.13$). These ratios compare to Tune~v1;
\S\ref{sec:3d-2p2h} uses untuned GENIE. Recovery as a by-product
of the higher-dimensional unfold provides the 4D$\to$3D analog of the
3D$\to$2D Jacobian anchor (\S\ref{sec:3d-anchor}).

\paragraph{What $q_{3}$ adds.} The resolved $(\Eavail,q_{3})$ plane
(Fig.~\ref{fig:q3excess}, left) reveals structure beneath the mild
$q_{3}$-marginal discrepancy integrated over $\Eavail$
($1.03$--$1.19$, center). At fixed $\Eavail$
in the quasielastic--$\Delta$ dip, data/Tune~v1 \emph{rises} with $q_{3}$.
The central value over-predicts low-$q_{3}$, moderate-$\Eavail$ cells
(ratio down to $0.66$); data excess concentrates at low $\Eavail$ (all $q_{3}$)
and toward high $q_{3}$. The muon-only or fixed-$q_{3}$ projections considered
here do not resolve this correlation, a natural input to the bin-compatible
Ascencio comparison. The map uses central values and the MINERvA Tune~v1
prior, distinct from the untuned GENIE sample of \S\ref{sec:3d-2p2h}
(App.~\ref{app:provenance}, \texttt{q3-excess}). No per-cell significances
are reported without a corrected 4D covariance (below).
% comparator weights: \texttt{w\_truth}; script \texttt{nd-unfolding/q3\_excess\_projection.py}

\paragraph{Why no 4D uncertainty is reported.} Corrected 4D covariance would use actual
minus/plus endpoints, MAT mean-centered $1/N$ blocks, one fixed estimator seed,
mean-centered joint throws with a separately reported mean shift, and no scalar
jitter subtraction. Because dump-all lateral branches are CV-support-limited,
selection-complete active-universe event loops must bound migration effects.
No 4D uncertainty magnitude or unified/block ratio is reported.

\paragraph{Bin-identical comparison to the published low-recoil measurement.}
Supplemental data from Ref.~\cite{MINERvA:2022incl} (44-cell
$d^{2}\sigma/(d\Eavail\,dq_{3})$ with its full covariance, distributed inside
the public arXiv source package) permit a direct quantitative cross-check.
Merging both measurements onto their maximal common grid --- edges present in
\emph{both} binnings --- and restricting our
$(\pT,\pPar)$ marginalization to $\pPar<\SI{20}{GeV/c}$ to mirror the Ascencio
muon gate ($\theta_\mu<20^{\circ}$, $1.5<p_\mu<\SI{20}{GeV}$) gives two super-cells,
$\Eavail<\SI{0.4}{GeV}$ in $q_{3}\in[0.4,0.6)$ and $[0.6,1.2)\,\si{GeV}$: the
low-recoil 2p2h region. Our central result lies
\SI{9.2}{\percent} and \SI{6.3}{\percent} above Ascencio. No pulls or
full-covariance $\chi^{2}$ are reported (no corrected 4D covariance, above;
Fig.~\ref{fig:ascencio}). Both share MINERvA flux/detector systematics,
treated here as independent. The $p_\mu$-vs-$\pPar$ gate differs by
$\lesssim\SI{6}{\percent}$ kinematic smearing at the \SI{20}{GeV} boundary.
A finer 44-cell comparison remains future work; it requires re-unfolding on
Ascencio edges (trivial for the unbinned method) \emph{and} a systematic
sweep on that binning.

\begin{figure}[H]
  \centering
  \includegraphics[width=0.85\textwidth]{ascencio_fullcov_compare_central}
  \caption{Bin-identical comparison to Ascencio
  \emph{et al.}~\cite{MINERvA:2022incl} on the maximal common
  $(\Eavail,q_{3})$ grid, used only for central-value comparison. Our points
  have no uncertainty or $\chi^{2}$ drawn because the
  quotable 4D covariance is not yet adopted (App.~\ref{app:history}); the
  published points carry the published covariance.}
  \label{fig:ascencio}
\end{figure}

\begin{figure}[H]
  \centering
  \includegraphics[width=0.99\textwidth]{q3_excess_projection}
  \caption{4D data/Tune~v1 ratio (``data/GENIE'' in the figure) projected onto
  $(\Eavail,q_{3})$ (left; blank cells are kinematically empty), with
  $q_{3}$ (center) and $\Eavail$ (right) marginals. The $\Eavail$ marginal
  reproduces the 3D low-recoil excess. The joint plane shows genuine $q_{3}$
  structure in the data\,$-$\,MC discrepancy (ratio rising with $q_{3}$
  across the quasielastic--$\Delta$ dip).}
  % Comparator: the MINERvA Tune~v1 prior (\texttt{w\_truth}).
  \label{fig:q3excess}
\end{figure}
